\documentclass[10pt,a4paper]{amsart}
\usepackage{amsmath,amssymb,amsthm,graphicx}
% Retained minimal preamble for the Stacks Project chapter more-algebra.tex (More on Algebra).
% The chapter file pinned-original-source.tex carries only \input{preamble} on its first line, so no
% preamble text lives inside the chapter itself.  The Stacks Project uses one single shared preamble for
% every chapter; the snapshot stages it identically as sources/stacks-chow/preamble.tex and
% sources/stacks-etale-cohomology/preamble.tex.  The block below is copied mechanically, line for line,
% from sources/stacks-chow/preamble.tex lines 14-18 (the Xy-pic setup) and 154-217 (the native theorem
% declarations together with the mathematical macros).  Nothing was added, removed or redefined; the
% wrapper class and the amsmath/amssymb/amsthm/graphicx base packages are supplied by the assembler.
% Chapter: more-algebra.tex -- More on Algebra.
% For commutative diagrams we use Xy-pic
\usepackage[all]{xy}

% We use 2cell for 2-commutative diagrams.
\xyoption{2cell}

% Theorem environments.
%
\theoremstyle{plain}
\newtheorem{theorem}[subsection]{Theorem}
\newtheorem{proposition}[subsection]{Proposition}
\newtheorem{lemma}[subsection]{Lemma}

\theoremstyle{definition}
\newtheorem{definition}[subsection]{Definition}
\newtheorem{example}[subsection]{Example}
\newtheorem{exercise}[subsection]{Exercise}
\newtheorem{situation}[subsection]{Situation}

\theoremstyle{remark}
\newtheorem{remark}[subsection]{Remark}
\newtheorem{remarks}[subsection]{Remarks}

\numberwithin{equation}{subsection}

% Macros
%
\def\lim{\mathop{\mathrm{lim}}\nolimits}
\def\colim{\mathop{\mathrm{colim}}\nolimits}
\def\Spec{\mathop{\mathrm{Spec}}}
\def\Hom{\mathop{\mathrm{Hom}}\nolimits}
\def\Ext{\mathop{\mathrm{Ext}}\nolimits}
\def\SheafHom{\mathop{\mathcal{H}\!\mathit{om}}\nolimits}
\def\SheafExt{\mathop{\mathcal{E}\!\mathit{xt}}\nolimits}
\def\Sch{\mathit{Sch}}
\def\Mor{\mathop{\mathrm{Mor}}\nolimits}
\def\Ob{\mathop{\mathrm{Ob}}\nolimits}
\def\Sh{\mathop{\mathit{Sh}}\nolimits}
\def\NL{\mathop{N\!L}\nolimits}
\def\CH{\mathop{\mathrm{CH}}\nolimits}
\def\proetale{{pro\text{-}\acute{e}tale}}
\def\etale{{\acute{e}tale}}
\def\QCoh{\mathit{QCoh}}
\def\Ker{\mathop{\mathrm{Ker}}}
\def\Im{\mathop{\mathrm{Im}}}
\def\Coker{\mathop{\mathrm{Coker}}}
\def\Coim{\mathop{\mathrm{Coim}}}

% Boxtimes
%
\DeclareMathSymbol{\boxtimes}{\mathbin}{AMSa}{"02}

%
% Macros for moduli stacks/spaces
%
\def\QCohstack{\mathcal{QC}\!\mathit{oh}}
\def\Cohstack{\mathcal{C}\!\mathit{oh}}
\def\Spacesstack{\mathcal{S}\!\mathit{paces}}
\def\Quotfunctor{\mathrm{Quot}}
\def\Hilbfunctor{\mathrm{Hilb}}
\def\Curvesstack{\mathcal{C}\!\mathit{urves}}
\def\Polarizedstack{\mathcal{P}\!\mathit{olarized}}
\def\Complexesstack{\mathcal{C}\!\mathit{omplexes}}
% \Pic is the operator that assigns to X its picard group, usage \Pic(X)
% \Picardstack_{X/B} denotes the Picard stack of X over B
% \Picardfunctor_{X/B} denotes the Picard functor of X over B
\def\Pic{\mathop{\mathrm{Pic}}\nolimits}
\def\Picardstack{\mathcal{P}\!\mathit{ic}}
\def\Picardfunctor{\mathrm{Pic}}
\def\Deformationcategory{\mathcal{D}\!\mathit{ef}}
\title{Lifting a bounded above acyclic complex of modules of the form P over IP to an acyclic complex whose terms lie in a class stable under summands and surjectivity modulo the ideal}
\author{The Stacks Project Authors}
\date{stacks-project:more-algebra.tex:local-snapshot-20261012; GNU FDL 1.2 or later}
\begin{document}
\maketitle
\setcounter{subsection}{1}
\noindent\textit{Editorial note.} All delivered mathematical text is the author's own, copied from the pinned Stacks Project chapter file more-algebra.tex; no statement or proof step was rewritten, reordered or removed. The counted statement is the labeled lemma environment at source lines 19342-19357 and the counted proof is the proof environment at source lines 19359-19387, which immediately follows it and is its own proof.\par\medskip
\noindent\textit{Editorial note.} The proof displays its induction diagram with Xy-pic. The retained preamble carries the Xy-pic setup of the shared Stacks preamble, so the diagram is compiled from the author's own markup; nothing about the diagram was changed.\par\medskip
\noindent\textit{Editorial note.} The chapter file pinned-original-source.tex carries only an input of the shared preamble on its first line, so no preamble text lives inside the chapter. The retained file sources/stacks-more-algebra/retained-preamble.tex is a mechanical, line-for-line copy of the project's single shared preamble as staged in the snapshot (Xy-pic setup plus the native theorem declarations and mathematical macros, including Ker, Im, colim and the bold L of the derived tensor product). Nothing was added, removed or redefined, and the excerpt serves only as preamble, never as a span coordinate.\par\medskip
\noindent\textit{Editorial note.} The source URL is the source-level Stacks Project URL, as required when more than one card is produced from one source; the per-claim tag is carried by claim\_key. The counted statement and proof refer to no labelled item of the chapter and cite no bibliography entry, so this card contributes no external reference stub and needs no bibliography.\par\medskip
\noindent\textit{Editorial note.} The section lead-in sentence at source lines 19337-19340, which poses the general lifting question and announces that the acyclic case is discussed first, is left outside the counted spans because it is an announcement and not part of the lemma or its proof.\par\medskip
\noindent\textit{Allowed standard prerequisites (editorial).}\begin{itemize}
\item A bounded above acyclic complex of modules with a chosen isomorphism P\textasciicircum{}i equal to Z\textasciicircum{}i direct sum Z\textasciicircum{}\{i+1\} has differentials given by the projection Z\textasciicircum{}i direct sum Z\textasciicircum{}\{i+1\} to Z\textasciicircum{}\{i+1\} followed by the inclusion of Z\textasciicircum{}\{i+1\} into Z\textasciicircum{}\{i+1\} direct sum Z\textasciicircum{}\{i+2\}. recorded in the proof at source lines 19373-19376 as the form of the maps P\textasciicircum{}i -> P\textasciicircum{}\{i+1\}. Turns the acyclicity of the lifted complex and the vanishing of its cohomology into the bookkeeping of the modules Z\textasciicircum{}i, which is what property (2) of the statement is applied to.
\item The lifting property of a projective module: a map into a quotient that factors through a surjection can be lifted, and the lift can be arranged to hit a prescribed element modulo the kernel. used at proof lines 19378-19382, where P\textasciicircum{}\{n-1\} is chosen projective and the composite P\textasciicircum{}\{n-1\} -> E\textasciicircum{}\{n-1\} -> E\textasciicircum{}n is lifted to a map P\textasciicircum{}\{n-1\} -> Z\textasciicircum{}n. Produces the new term P\textasciicircum{}\{n-1\} together with the differential into Z\textasciicircum{}n that extends the induction diagram one step to the left.
\item The quotient map P -> P/IP is compatible with the given isomorphism P\textasciicircum{}i/IP\textasciicircum{}i -> E\textasciicircum{}i, so that the identification Z\textasciicircum{}n/IZ\textasciicircum{}n = Im(E\textasciicircum{}\{n-1\} -> E\textasciicircum{}n) holds. used at proof lines 19372 and 19380. Identifies the target of the lifted map with the quotient of Z\textasciicircum{}n, which is exactly the condition needed for the lifting property to apply.
\end{itemize}\medskip
\begin{lemma}
\label{lemma-lift-acyclic-complex}
Let $R$ be a ring. Let $I \subset R$ be an ideal. Let $\mathcal{P}$
be a class of $R$-modules. Assume
\begin{enumerate}
\item each $P \in \mathcal{P}$ is a projective $R$-module,
\item if $P_1 \in \mathcal{P}$ and $P_1 \oplus P_2 \in \mathcal{P}$, then
$P_2 \in \mathcal{P}$, and
\item if $f : P_1 \to P_2$, $P_1, P_2 \in \mathcal{P}$ is surjective
modulo $I$, then $f$ is surjective.
\end{enumerate}
Then given any bounded above acyclic complex $E^\bullet$
whose terms are of the form $P/IP$ for $P \in \mathcal{P}$ there
exists a bounded above acyclic complex $P^\bullet$ whose terms
are in $\mathcal{P}$ lifting $E^\bullet$.
\end{lemma}

\begin{proof}
Say $E^i = 0$ for $i > b$. Assume given $n$ and a morphism of complexes
$$
\xymatrix{
& & P^n \ar[r] \ar[d] & P^{n + 1} \ar[r] \ar[d] & \ldots \ar[r]
& P^b \ar[r] \ar[d] & 0 \ar[r] \ar[d] & \ldots \\
\ldots \ar[r] & E^{n - 1} \ar[r] &
E^n \ar[r] & E^{n + 1} \ar[r] & \ldots \ar[r] &
E^b \ar[r] & 0 \ar[r] & \ldots
}
$$
with $P^i \in \mathcal{P}$, with
$P^n \to P^{n + 1} \to \ldots \to P^b$ acyclic in degrees $\geq n + 1$,
and with vertical maps inducing isomorphisms $P^i/IP^i \to E^i$.
In this situation one can inductively choose isomorphisms
$P^i = Z^i \oplus Z^{i + 1}$ such that the maps $P^i \to P^{i + 1}$
are given by
$Z^i \oplus Z^{i + 1} \to Z^{i + 1} \to Z^{i + 1} \oplus Z^{i + 2}$.
By property (2) and arguing inductively we see that $Z^i \in \mathcal{P}$.
Choose $P^{n - 1} \in \mathcal{P}$ and an isomorphism
$P^{n - 1}/IP^{n - 1} \to E^{n - 1}$. Since
$P^{n - 1}$ is projective and since $Z^n/IZ^n = \Im(E^{n - 1} \to E^n)$,
we can lift the map $P^{n - 1} \to E^{n - 1} \to E^n$ to a map
$P^{n - 1} \to Z^n$. By property (3) the map $P^{n - 1} \to Z^n$
is surjective. Thus we obtain an extension of the diagram
by adding $P^{n - 1}$ and the maps just constructed to the left
of $P^n$. Since a diagram of the desired form exists for $n > b$
we conclude by induction on $n$.
\end{proof}

\end{document}
